% Recovered from the previously completed PDF; no new mathematical content.
% Prose is editable; formula boxes retain original symbol positions.
% Build from repository root: xelatex -no-shell-escape -halt-on-error -output-directory=... print/fudan/NAME.tex
\documentclass[11pt,a4paper]{article}
\usepackage[margin=23mm,headheight=14pt,headsep=8.5mm,footskip=10mm]{geometry}
\usepackage{amsmath,amssymb}
\DeclareMathAlphabet{\mathsf}{OT1}{cmss}{m}{n}
\usepackage[no-math]{fontspec}
\usepackage{xeCJK}
\setmainfont{Latin Modern Roman}
\setCJKmainfont{FandolSong-Regular.otf}[Path=assets/fonts/fandol/,BoldFont=FandolSong-Bold.otf,ItalicFont=FandolSong-Regular.otf,Script=CustomDefault]
\xeCJKsetup{PunctStyle=quanjiao,CheckSingle=true}
\usepackage{setspace}\setstretch{1.12}
\setlength{\parindent}{2em}\setlength{\parskip}{3pt}
\emergencystretch=6em
\hfuzz=9pt
\usepackage{titlesec}
\titleformat{\section}{\centering\normalfont\Large\bfseries}{}{0pt}{}
\titlespacing*{\section}{0pt}{1.8\baselineskip}{1.0\baselineskip}
\usepackage{fancyhdr}\pagestyle{fancy}\fancyhf{}
\fancyhead[L]{\small 数学手记·高等代数}\fancyhead[R]{\small 20S14–20S18、21A01–21A02}
\fancyfoot[C]{\thepage}\renewcommand{\headrulewidth}{0.4pt}
\usepackage[hidelinks,unicode]{hyperref}
\hypersetup{pdftitle={20S14–20S18、21A01–21A02：题目、解答与小结},pdfauthor={Vincent Arqain}}
\newcommand{\RB}[4]{\raisebox{0pt}[#2bp][#3bp]{\hbox to #1bp{#4\hss}}}
\newcommand{\RG}[5]{\hbox to 0pt{\kern#1bp\raisebox{#2bp}[0pt][0pt]{\fontsize{#3bp}{#3bp}\csname font#4\endcsname\char#5}\hss}}
\newcommand{\RT}[4]{\hbox to 0pt{\kern#1bp\raisebox{#2bp}[0pt][0pt]{\fontsize{#3bp}{#3bp}\selectfont#4}\hss}}
\newcommand{\RS}[4]{\hbox to 0pt{\kern#1bp\raisebox{#2bp}[0pt][0pt]{\fontsize{#3bp}{#3bp}\selectfont\ensuremath{#4}}\hss}}
\newcommand{\RR}[4]{\hbox to 0pt{\kern#1bp\raisebox{#2bp}[0pt][0pt]{\rule{#3bp}{#4bp}}\hss}}
\newcommand{\fontmi}{\usefont{OML}{cmm}{m}{it}}
\newcommand{\fontsy}{\usefont{OMS}{cmsy}{m}{n}}
\newcommand{\fontex}{\usefont{OMX}{cmex}{m}{n}}
\newcommand{\fontrm}{\usefont{OT1}{cmr}{m}{n}}
\newcommand{\fontbf}{\usefont{OT1}{cmr}{bx}{n}}
\newcommand{\fontam}{\usefont{U}{msa}{m}{n}}
\newcommand{\fontbb}{\usefont{U}{msb}{m}{n}}
\widowpenalty=10000\clubpenalty=10000\raggedbottom
\DeclareMathOperator{\rank}{rank}
\DeclareMathOperator{\tr}{tr}
\DeclareMathOperator{\diag}{diag}
\DeclareMathOperator{\adj}{adj}
\DeclareMathOperator{\Ker}{Ker}
\DeclareMathOperator{\Imop}{Im}
\newcounter{fudanquestion}
\newcommand{\question}[1]{\stepcounter{fudanquestion}\section*{#1}}
\newcommand{\sol}{\par\noindent 解：\ }
\newcommand{\method}{\par\noindent 方法提炼：\ }
\newcommand{\sourcepage}[1]{\par\noindent{\small 题目来源：谢启鸿《复旦大学高等代数习题集（2023 版）》，第 #1 页。}}
\newcommand{\proofend}{\hfill$\square$\par}
\newcommand{\id}{I_n}

\begin{document}
\begin{center}{\LARGE\bfseries 20S14–20S18、21A01–21A02\par}\vspace{8pt}复旦大学高等代数习题集·题目与解答\end{center}


\question{2020S14}
\noindent 设 $n$ 元实二次型
\[
  f(x_1,\ldots,x_n)=\sum_{i,j=1}^{n}\max\{i,j\}\,x_i x_j.
\]
求此二次型的规范标准型。

\sol 记 $s=\sum_{i=1}^{n}x_i$。本题的关键是将系数 $\max\{i,j\}$ 改写为与部分和有关的表达式。对任意 $1\leq i,j\leq n$，有
\[
\max\{i,j\}=n-\sum_{k=1}^{n-1}
 \mathbf 1_{\{i\leq k\}}\mathbf 1_{\{j\leq k\}},
\]
其中 $\mathbf 1_E$ 表示事件 $E$ 的示性函数。事实上，上式求和项为满足 $k\geq\max\{i,j\}$ 的整数 $k\in\{1,\ldots,n-1\}$ 的个数，恰好等于 $n-\max\{i,j\}$。

因此，交换有限求和的次序后，
\begin{align*}
 f(x_1,\ldots,x_n)
 &=n\Bigl(\sum_{i=1}^{n}x_i\Bigr)^2
   -\sum_{k=1}^{n-1}
    \sum_{i,j=1}^{k}x_i x_j\\
 &=n\Bigl(\sum_{i=1}^{n}x_i\Bigr)^2
   -\sum_{k=1}^{n-1}\Bigl(\sum_{i=1}^{k}x_i\Bigr)^2.
\end{align*}
作可逆线性换元
\[
 y_k=x_1+\cdots+x_k\quad(1\leq k\leq n),
 \qquad x_1=y_1,\quad x_k=y_k-y_{k-1}\ (k\geq2).
\]
便得到合同对角形式
\[
 f=n y_n^2-y_1^2-y_2^2-\cdots-y_{n-1}^2.
\]
再令 $z_n=\sqrt n\,y_n$，$z_k=y_k$（$1\leq k\leq n-1$），可得规范标准型
\[
 {\ f=z_n^2-z_1^2-\cdots-z_{n-1}^2\ }.
\]
因此，该二次型的正惯性指数为 $1$，负惯性指数为 $n-1$，秩为 $n$。特别地，相应对称矩阵
$M=(\max\{i,j\})_{n\times n}$ 满足
\[
 \det M=(-1)^{n-1}n.
\]
这是因为 $x\mapsto y$ 的变换矩阵是行列式为 $1$ 的下三角阵。
\proofend
\method 对“最大值”型矩阵，可用{前缀和换元}。恒等式
$\max\{i,j\}=n-(n-\max\{i,j\})$ 把原二次型化为一个正平方项与 $n-1$ 个负平方项之和，无须计算特征值。
\sourcepage{94}

\question{2020S15}
\noindent 设 $M$ 为 $n$ 阶实方阵。若对任意 $x\in\mathbb R^n$，均有
\(x^{\mathsf T}Mx\geq0\)，则称 $M$ 为{亚半正定阵}。

\par\smallskip\noindent{(1)} 证明下列条件等价：$M$ 为亚半正定阵；$M+M^{\mathsf T}$ 为半正定阵；存在半正定实对称阵 $H$ 及实反对称阵 $S$，使得 $M=H+S$。

\par\smallskip\noindent{(2)} 设 $A,B$ 为亚半正定阵，$c\geq0$。证明：
\par\noindent\quad (2.1) $A+B$、$cA$、$A^{\mathsf T}$ 和伴随阵 $A^*$ 都是亚半正定阵；
\par\noindent\quad (2.2) 对任意实方阵 $C$，$C^{\mathsf T}AC$ 也是亚半正定阵；
\par\noindent\quad (2.3) 若 $C$ 为亚正定阵，则 $A+C$ 也是亚正定阵；
\par\noindent\quad (2.4) $A$ 的所有特征值的实部均非负，且 $\det A\geq0$；
\par\noindent\quad (2.5) 举例说明：非异亚半正定阵不一定是亚正定阵。

\sol{(1)} 对任意实列向量 $x$，标量 $x^{\mathsf T}Mx$ 等于其转置，故
\[
 x^{\mathsf T}Mx=x^{\mathsf T}M^{\mathsf T}x
   =\frac12 x^{\mathsf T}(M+M^{\mathsf T})x.
\]
于是 $M$ 亚半正定当且仅当对称阵 $M+M^{\mathsf T}$ 半正定。另一方面，任意实方阵都有唯一分解
\[
 M=\underbrace{\frac{M+M^{\mathsf T}}2}_{H}
   +\underbrace{\frac{M-M^{\mathsf T}}2}_{S},
 \qquad H^{\mathsf T}=H,\quad S^{\mathsf T}=-S.
\]
由于 $x^{\mathsf T}Sx=0$，故 $M$ 亚半正定当且仅当 $H$ 半正定，三个条件因而等价。

\par\smallskip\noindent{(2.1) 前三类矩阵。} 对任意 $x\in\mathbb R^n$，由已知条件直接得到
\[
\begin{aligned}
 x^{\mathsf T}(A+B)x&=x^{\mathsf T}Ax+x^{\mathsf T}Bx\geq0,\\
 x^{\mathsf T}(cA)x&=c\,x^{\mathsf T}Ax\geq0,\\
 x^{\mathsf T}A^{\mathsf T}x&=x^{\mathsf T}Ax\geq0.
\end{aligned}
\]
因此三者均亚半正定。关于伴随阵 $A^*$ 的断言稍后由谱性质及摄动法证明。

\par\smallskip\noindent{(2.2)} 对任意 $x\in\mathbb R^n$，有
\[
 x^{\mathsf T}(C^{\mathsf T}AC)x
  =(Cx)^{\mathsf T}A(Cx)\geq0.
\]
故 $C^{\mathsf T}AC$ 亚半正定；这里不要求 $C$ 可逆。

\par\smallskip\noindent{(2.3)} 由亚正定阵的定义，对任何 $x\neq0$，有
\[
 x^{\mathsf T}(A+C)x=x^{\mathsf T}Ax+x^{\mathsf T}Cx
  \geq x^{\mathsf T}Cx>0.
\]
因此 $A+C$ 亚正定。

\par\smallskip\noindent{(2.4)} 设 $Av=\lambda v$，其中 $0\neq v\in\mathbb C^n$。由 (1) 知 $H=(A+A^{\mathsf T})/2$ 半正定，因此
\[
 2\operatorname{Re}\lambda\,\lVert v\rVert_2^2
 =v^*(A+A^{\mathsf T})v
 =2v^*Hv\geq0.
\]
此处 $v^*$ 表示共轭转置。因此 $\operatorname{Re}\lambda\geq0$。由于 $A$ 是实矩阵，非实特征值成共轭对出现，每一对的乘积为 $|\lambda|^2\geq0$；其余实特征值均非负。于是所有特征值之积满足
\[
 {\det A\geq0}.
\]
此外，若 $A$ 亚正定，则上面的不等号严格成立，从而每个特征值的实部都严格为正，特别有 $\det A>0$。

\par\smallskip\noindent{(2.1) 伴随阵。} 此处按原书记号，$A^*=\adj(A)$ 表示{经典伴随矩阵}，并非共轭转置。取任意 $\varepsilon>0$，令
\[
 A_\varepsilon=A+\varepsilon I_n.
\]
由 (2.3)，$A_\varepsilon$ 亚正定，因而由刚证得的谱性质，$\det A_\varepsilon>0$ 且 $A_\varepsilon$ 可逆。对任意 $x\neq0$，令 $y=A_\varepsilon^{-1}x\neq0$，则
\[
 x^{\mathsf T}A_\varepsilon^{-1}x
 =(A_\varepsilon y)^{\mathsf T}y
 =y^{\mathsf T}A_\varepsilon y>0.
\]
故 $A_\varepsilon^{-1}$ 亚正定，进而
\[
 x^{\mathsf T}\adj(A_\varepsilon)x
 =\det(A_\varepsilon)\,x^{\mathsf T}A_\varepsilon^{-1}x\geq0.
\]
伴随阵各元素是矩阵元素的多项式，所以 $\adj(A_\varepsilon)\to\adj(A)$。令 $\varepsilon\downarrow0$，得到
\[
 x^{\mathsf T}A^*x\geq0\qquad(x\in\mathbb R^n),
\]
即 $A^*$ 也是亚半正定阵。

\par\smallskip\noindent{(2.5)} 取二阶实反对称阵
\[
 A=\begin{pmatrix}0&-1\\1&0\end{pmatrix}.
\]
此时 $\det A=1$，故 $A$ 非异；但对任意 $x\in\mathbb R^2$，恒有 $x^{\mathsf T}Ax=0$。因此 $A$ 亚半正定，却不是亚正定阵。
\proofend
\method 所有结论统一于{对称部分} $H=(A+A^{\mathsf T})/2$；谱的实部由 $v^*Hv$ 控制。非异情形的伴随阵用 $\adj(A)=\det(A)A^{-1}$ 处理，奇异情形再借 $A+\varepsilon I_n$ 的摄动取极限。
\sourcepage{94}

\question{2020S16}
\noindent 设 $a>0$。证明下列 $n$ 阶实对称矩阵正定：
\[
 A=\begin{pmatrix}
 a & \dfrac{a^2}{2} & \cdots & \dfrac{a^n}{n}\\[4pt]
 \dfrac{a^2}{2} & \dfrac{a^3}{3} & \cdots & \dfrac{a^{n+1}}{n+1}\\
 \vdots & \vdots & \ddots & \vdots\\
 \dfrac{a^n}{n} & \dfrac{a^{n+1}}{n+1} & \cdots & \dfrac{a^{2n-1}}{2n-1}
 \end{pmatrix}.
\]

\sol 观察到 $A$ 的第 $(i,j)$ 个元素为
\[
 a_{ij}=\frac{a^{i+j-1}}{i+j-1}
         =\int_0^a t^{i+j-2}\,\mathrm dt
         \qquad(1\leq i,j\leq n).
\]
任取非零实列向量 $x=(x_1,\ldots,x_n)^{\mathsf T}$，令
\[
 p(t)=x_1+x_2t+\cdots+x_nt^{n-1}.
\]
那么，利用有限求和与积分的交换，
\begin{align*}
 x^{\mathsf T}Ax
 &=\sum_{i,j=1}^{n}x_i a_{ij}x_j
  =\sum_{i,j=1}^{n}x_ix_j\int_0^a t^{i+j-2}\,\mathrm dt\\
 &=\int_0^a\left(\sum_{i=1}^{n}x_i t^{i-1}\right)^2\mathrm dt
  =\int_0^a p(t)^2\,\mathrm dt.
\end{align*}
由于 $x\neq0$，$p$ 是非零多项式。非零多项式不可能在整个区间 $[0,a]$ 上恒为零；又 $p(t)^2$ 连续且非负，因此
\[
 \int_0^a p(t)^2\,\mathrm dt>0.
\]
故对任意 $x\neq0$ 都有 $x^{\mathsf T}Ax>0$，即 $A$ 是正定实对称阵。
\proofend

\method $A$ 实际上是 $L^2[0,a]$ 中线性无关函数组 $1,t,\ldots,t^{n-1}$ 的{Gram 矩阵}。对矩阵元素形如 $\int\!\phi_i(t)\phi_j(t)\,\mathrm dt$ 的题目，优先把二次型写成一个平方积分；正定性便转化为函数组的线性无关性。

\noindent{补充。} 令 $H_n=(1/(i+j-1))_{n\times n}$，$D=\diag(a^{1/2},a^{3/2},\ldots,a^{n-1/2})$，则
\[
 A=DH_nD,\qquad \det A=a^{n^2}\det H_n>0.
\]
这说明本题是 Hilbert 矩阵正定性的一个带尺度参数的版本。
\sourcepage{95}

\question{2020S17}
\noindent 设 $V$ 为区间 $[-1,1]$ 上次数不超过 $3$ 的实系数多项式构成的线性空间，内积定义为
$\langle f,h\rangle=\int_{-1}^{1}f(x)h(x)\,\mathrm dx$。求
\[
 \min_{f\in V}\int_{-1}^{1}\bigl(\sin(\pi x)-f(x)\bigr)^2\,\mathrm dx.
\]

\sol 令 $g(x)=\sin(\pi x)$。所求即 $g$ 到有限维内积子空间 $V$ 的{最佳平方逼近}。由于 $g$ 是奇函数，且奇函数与偶函数在对称区间上的内积为零，故只需考虑 $V$ 的奇多项式部分 $L(x,x^3)$。

该二维空间的一组标准正交基是归一化的 Legendre 多项式
\[
 w_1(x)=\sqrt{\frac32}\,x,\qquad
 w_3(x)=\sqrt{\frac78}\,(5x^3-3x).
\]
由分部积分可得
\[
 \int_{-1}^{1}x\sin(\pi x)\,\mathrm dx=\frac2\pi,
 \qquad
 \int_{-1}^{1}x^3\sin(\pi x)\,\mathrm dx
 =\frac2\pi-\frac{12}{\pi^3}.
\]
因此，$g$ 在这两个正交方向上的投影系数分别为
\[
\begin{aligned}
 \langle g,w_1\rangle&=\sqrt{\frac32}\cdot\frac2\pi
     =\frac{\sqrt6}\pi,\\
 \langle g,w_3\rangle&=\sqrt{\frac78}\left[5\left(\frac2\pi-\frac{12}{\pi^3}\right)-\frac6\pi\right]
     =\frac{\sqrt{14}(\pi^2-15)}{\pi^3}.
\end{aligned}
\]
由正交投影定理，唯一的最佳逼近多项式为
\begin{align*}
 f_*(x)&=\langle g,w_1\rangle w_1(x)
       +\langle g,w_3\rangle w_3(x)\\
 &=\frac{15(21-\pi^2)}{2\pi^3}x
   +\frac{35(\pi^2-15)}{2\pi^3}x^3.
\end{align*}
同时，$\lVert g\rVert^2=\int_{-1}^{1}\sin^2(\pi x)\,\mathrm dx=1$。利用勾股恒等式，得到
\begin{align*}
 \min_{f\in V}\lVert g-f\rVert^2
 &=\lVert g\rVert^2-|\langle g,w_1\rangle|^2
                         -|\langle g,w_3\rangle|^2\\
 &={\ 1-\frac6{\pi^2}-\frac{14(\pi^2-15)^2}{\pi^6}\ }\\
 &\approx0.00878023323873.
\end{align*}
等号当且仅当 $f=f_*$ 时成立，故上式就是所求最小值。
\proofend
\method 本题的核心是{正交投影与奇偶性}。先把四维多项式空间缩减为二维奇多项式空间，再用正交 Legendre 基计算两个内积；这样避免了对四个待定系数建立正规方程。
\sourcepage{95}

\question{2020S18}
\noindent 设 $Q$ 为 $n$ 阶实正交矩阵，且 $1$ 不是 $Q$ 的特征值。设
\[
 P=I_n-2uu^{\mathsf T},\qquad u\in\mathbb R^n,\quad u^{\mathsf T}u=1.
\]
证明：$1$ 是 $PQ$ 的特征值。

\sol 因 $1$ 不是 $Q$ 的特征值，$I_n-Q$ 可逆。构造非零列向量
\[
 v=(I_n-Q)^{-1}u.
\]
于是
\[
 (I_n-Q)v=u,\qquad Qv=v-u.
\]
利用 $Q$ 的正交性，得到 $\lVert Qv\rVert_2=\lVert v\rVert_2$，故
\[
 \lVert v-u\rVert_2^2=\lVert v\rVert_2^2.
\]
展开左端，结合 $u^{\mathsf T}u=1$，有
\[
 \lVert v\rVert_2^2-2u^{\mathsf T}v+1
   =\lVert v\rVert_2^2,
 \qquad\text{从而}\quad u^{\mathsf T}v=\frac12.
\]
因此
\[
 u^{\mathsf T}Qv
 =u^{\mathsf T}(v-u)
 =\frac12-1=-\frac12.
\]
最后直接计算：
\begin{align*}
 PQv&=(I_n-2uu^{\mathsf T})Qv\\
    &=Qv-2u(u^{\mathsf T}Qv)\\
    &=(v-u)-2u\left(-\frac12\right)=v.
\end{align*}
由于 $u\neq0$ 且 $I_n-Q$ 可逆，故 $v=(I_n-Q)^{-1}u\neq0$。这样便构造出了 $PQ$ 的属于特征值 $1$ 的非零特征向量，结论得证。
\proofend
\method 反射矩阵 $P=I_n-2uu^{\mathsf T}$ 只沿 $u$ 的方向改变向量。由 $I_n-Q$ 可逆，主动寻找满足 $(I_n-Q)v=u$ 的向量，再利用 $\|Qv\|=\|v\|$，即可得到 $PQv=v$；这比直接展开 $\det(I_n-PQ)$ 更直观。
\sourcepage{95}

\question{2021A01}
\noindent 求下列 $n$ 阶行列式的值：
\[
 D_n=\begin{vmatrix}
 1&a&a^2&\cdots&a^{n-1}\\
 a&1&a&\cdots&a^{n-2}\\
 a^2&a&1&\cdots&a^{n-3}\\
 \vdots&\vdots&\vdots&\ddots&\vdots\\
 a^{n-1}&a^{n-2}&a^{n-3}&\cdots&1
 \end{vmatrix}.
\]

\sol 设 $A=(a^{|i-j|})_{1\leq i,j\leq n}$。对行列式依次施行初等行变换
\[
 R_i\longleftarrow R_i-aR_{i-1},
 \qquad i=n,n-1,\ldots,2.
\]
这里必须{从下往上}操作，以保证每一步用到的 $R_{i-1}$ 仍是原来的行。记变换后的第 $(i,j)$ 个元素为 $b_{ij}$。当 $i\geq2$ 时，
\[
 b_{ij}=a^{|i-j|}-a\,a^{|i-1-j|}
 =\begin{cases}
 0,&j<i,\\[2pt]
 1-a^2,&j=i,\\[2pt]
 (1-a^2)a^{j-i},&j>i.
 \end{cases}
\]
第一行不变。因此，所得矩阵为上三角阵，主对角元依次为
\[
 1,\quad 1-a^2,\quad\ldots,\quad1-a^2.
\]
第三类初等行变换不改变行列式的值，故
\[
 {\ D_n=(1-a^2)^{n-1}\ }.
\]
当 $n=1$ 时公式也成立；当 $a=\pm1$ 且 $n\geq2$ 时，行列式确为零。
\proofend
\method 本题属于{等比型对称 Toeplitz 行列式}。沿相邻两行作 $R_i-aR_{i-1}$，并注意自下而上的操作次序，可一次消去主对角线以下全部元素。对于实数 $a$ 且 $n\geq2$，还可由各阶顺序主子式立即判断：该矩阵正定当且仅当 $|a|<1$。
\sourcepage{107}

\question{2021A02}
\noindent 设 $n\geq2$，求行列式
\[
 D=\begin{vmatrix}
 1&x_1&\cdots&x_1^{n-2}&(x_2+\cdots+x_n)^n\\
 1&x_2&\cdots&x_2^{n-2}&(x_1+x_3+\cdots+x_n)^n\\
 \vdots&\vdots&\ddots&\vdots&\vdots\\
 1&x_n&\cdots&x_n^{n-2}&(x_1+\cdots+x_{n-1})^n
 \end{vmatrix}.
\]

\sol 记
\[
 s=x_1+\cdots+x_n,\qquad
 \Delta=\prod_{1\leq i<j\leq n}(x_j-x_i).
\]
原行列式第 $i$ 行末项为 $(s-x_i)^n$。由二项式定理，
\[
 (s-x_i)^n
 =\sum_{k=0}^{n}(-1)^k\binom nk s^{n-k}x_i^k.
\]
由于原行列式的前 $n-1$ 列依次是 $1,x_i,\ldots,x_i^{n-2}$，可用这些列消去最后一列中所有 $x_i^k$（$0\leq k\leq n-2$）的项，而行列式不变。故
\[
 D=(-1)^{n-1}ns\,\Delta+(-1)^n E,
\]
其中
\[
 E=\det\begin{pmatrix}
 1&x_1&\cdots&x_1^{n-2}&x_1^n\\
 1&x_2&\cdots&x_2^{n-2}&x_2^n\\
 \vdots&\vdots&\ddots&\vdots&\vdots\\
 1&x_n&\cdots&x_n^{n-2}&x_n^n
 \end{pmatrix}.
\]
为了计算 $E$，考虑首一多项式
\[
 p(t)=\prod_{j=1}^{n}(t-x_j)
     =t^n-st^{n-1}+q(t),\qquad\deg q\leq n-2.
\]
将 $t=x_i$ 代入 $p(t)=0$，便有
\[
 x_i^n=sx_i^{n-1}-q(x_i)\qquad(1\leq i\leq n).
\]
因此 $E$ 的最后一列可替换成 $sx_i^{n-1}$，因为 $q(x_i)$ 对应的列是前 $n-1$ 列的线性组合。根据 Vandermonde 行列式公式，
\[
 E=s\det(x_i^{j-1})_{1\leq i,j\leq n}=s\Delta.
\]
代回原式，得到
\begin{align*}
 D&=(-1)^{n-1}ns\Delta+(-1)^ns\Delta\\
  &=(-1)^{n-1}(n-1)s\Delta.
\end{align*}
即
\[
 {\ D=(-1)^{n-1}(n-1)
 \left(\sum_{i=1}^{n}x_i\right)
 \prod_{1\leq i<j\leq n}(x_j-x_i)\ }.
\]
上述推导完全是多项式恒等式，因此即使若干 $x_i$ 相同，结论仍然成立。
\proofend
\method 对于{Vandermonde 型行列式}，先利用已有的 $1,x,\ldots,x^{n-2}$ 列把高次多项式削减成两项。剩余的 $x^n$ 列利用根多项式 $\prod(t-x_i)$ 降为 $s x^{n-1}$，可避免繁琐的行列式展开。
\sourcepage{107}


\end{document}
